\documentclass[10pt]{article}
\usepackage[letterpaper,margin=0.68in]{geometry}
\usepackage[T1]{fontenc}
\usepackage{lmodern,microtype,amsmath,amssymb,booktabs,array,enumitem,xcolor,fancyhdr}
\usepackage[hidelinks]{hyperref}
\definecolor{ink}{HTML}{20354B}
\pagestyle{fancy}\fancyhf{}\fancyhead[L]{\small\color{ink}Model solution and calibration}\fancyhead[R]{\small September 2026}
\fancyfoot[C]{\small\thepage\ / 2}\renewcommand{\headrulewidth}{0.3pt}
\setlength{\headheight}{14pt}\setlength{\parindent}{0pt}\setlength{\parskip}{4.5pt}
\setlist[enumerate]{leftmargin=*,itemsep=5pt,topsep=3pt,parsep=2pt}
\newcommand{\lead}[1]{\textbf{\color{ink}#1}}
\begin{document}
\begin{center}
{\Large\bfseries The stationary solution: the actual algorithm}\\[3pt]
{\small Guess prices, solve households, aggregate, update prices, repeat}
\end{center}

\lead{What ``stationary'' means here.} Prices and aggregates are constant; households still age, save and have children. Hold economic parameters fixed. The current model has one market, four-year periods, 17 age cells and 120 wealth nodes. States also record inherited tenure/product, income, lifetime births and children at home.

\begin{enumerate}
\item \lead{Guess the house price $p$. Compute rent and household resources.}
With a constant price, no-arbitrage gives $r=(R-1+\delta+\tau_H)p$, where $R$ is the gross financial return, $\delta$ depreciation and $\tau_H$ the property-tax rate, all in four-year units. Given $p$ and $r$, construct each state's income, sale proceeds, housing costs, borrowing/down-payment limits and feasible choices. The retained historical baseline fixes transfers at zero.

\item \lead{Start at the oldest age; solve conditional household problems.}
The final age uses the terminal bequest payoff. Then move one age backward at a time. At each state, consider no birth and an additional birth, and each feasible rental/owner product. For each branch and product, choose saving and current consumption/housing using the already computed next-age values. Denote this optimized deterministic value by $Q_{d,k}$, where $d\in\{0,1\}$ is the birth outcome and $k$ the housing product.

For a renter at given saving, split resources net of family consumption/housing requirements using Cobb--Douglas shares; cap rental housing and put the remainder into consumption. For an owner product, housing is fixed and the budget gives consumption. Only a one-dimensional saving search remains. For saving, search every feasible piecewise-linear continuation segment: compare endpoints, rental-cap kinks and interior first-order-condition solutions. This is exhaustive for the numerical approximation. Income risk, survival, child departures and bequests enter continuation values. Store policies and values; \textbf{repeat for every state at that age.}

\item \lead{Combine these values in the sequential choice operator.}
First integrate the \emph{subsequent} housing shocks conditional on the birth outcome. The housing inclusive value is:
\[
 H_d=\kappa\log\sum_k\exp(Q_{d,k}/\kappa).
\]
Here $\kappa$ is the housing-choice taste scale. $H_d$ is the expected optimized value of housing and saving after outcome $d$. Now construct the values of waiting, $W$, and attempting, $A$, before biological success is known:
\[
 W=H_0,\qquad A=(1-\pi)H_0+\pi(H_1-C_1),\qquad
 P_{\rm attempt}=\frac{e^{A/\sigma}}{e^{W/\sigma}+e^{A/\sigma}}.
\]
The biological success probability is $\pi$. The utility cost $C_1$ applies only to a first birth; it is zero for later births. The fertility taste scale $\sigma$ can differ between first and later attempts. Birth probability is $\pi P_{\rm attempt}$. The pre-choice value passed to the preceding age is:
\[
 V=\sigma\log\{e^{W/\sigma}+e^{A/\sigma}\}.
\]
Economic timing is attempt, biological outcome, then housing/saving; computation reverses this order. Ineligible states use the no-attempt branch. \textbf{Return to step 2 for the next younger age.} All taste-shock integration is analytic.

\item \lead{Hold policies fixed; construct the stationary cross-section.}
In the old stationary economy, start with the entering cohort and propagate its mass forward through all ages using choice probabilities and income, birth, child-departure and survival transitions. Constant entry generates the cross-section. The distant endpoint additionally iterates its household/person demographic fixed point. Aggregate housing policies against household mass to obtain $D(p)$.

\item \lead{Compare demand with supply; change the price and repeat.}
Compute the signed relative market error $F(p)=[D(p)-S(p)]/S(p)$. If demand exceeds supply, search toward a higher price; if below, toward a lower one. \textbf{At each new price, return to step 1 and re-solve household policies and the distribution.}

The active old initializer brackets the scalar price root and refines it with Brent's method. The matched endpoint uses a safeguarded bracketed root in log price. A narrow bracket is insufficient: the actual market residual, accounting and feasibility checks must pass; repeat the final solution.
\end{enumerate}

\lead{What changes in the simultaneous arm?} Replace step 3 by a fertility-nested logsum over complete plans: one housing product for waiting; failure/success products for attempting. Shocks realize jointly before biological success; housing scale cannot exceed fertility scale. Household, distribution and price loops remain.

\newpage
\begin{center}
{\Large\bfseries The transition and calibration algorithms}\\[3pt]
{\small The stationary calculation is a building block inside two larger loops}
\end{center}

\lead{A. Solve the perfect-foresight transition at fixed parameters.}
Households learn the aggregate path in 2007; 2023 is a measurement date on that path. Individual outcomes remain uncertain. Immediate knowledge need not mean an immediate full preference change.

\begin{enumerate}
\item \textbf{Construct the initial state and solve a stationary endpoint.} Use the old stationary economy for inherited states and pre-2007 cohort history, with observed age/head margins through 2023. Solve the distant stationary economy under terminal preferences, demographics and fiscal rules for $V_{T+1}$ and $p_{T+1}$. This endpoint is fixed during the price-path iteration.
\item \textbf{Guess the entire price vector $(p_{2007},p_{2011},\ldots,p_T)$.} Here $t+1$ denotes the next four-year model date. Compute rent from:
\[
 r_t=(R+\delta+\tau_H)p_t-p_{t+1}.
\]
Capital gains affect current rents and choices. Adjust infeasible-rent guesses. The terminal date is numerical closure, not an empirical target year.
\item \textbf{Work backward in calendar time.} At $T$, solve households as on page 1 using next date's values. Move to the preceding model date and repeat back to 2007, storing dated values. All decisions use the same anticipated path.
\item \textbf{Start again in 2007 and work forward.} Apply dated policies to inherited mass, measure demand and advance the distribution. Use historical demographic conditioning through 2023 and the person/head law afterward. Never reset each date to a stationary distribution. Recomputing policies from stored values gives two household-solution passes per date.
\item \textbf{Update the whole price vector and return to step 2.} Compute $F_t=(D_t-S_t)/S_t$ at every date. A safeguarded Broyden method updates all log prices jointly, initialized from finite-difference price responses. Each trial repeats the full backward/forward evaluation. Require all market errors below $2\times10^{-4}$ (0.02\%), other gates and fresh reproduction.
\item \textbf{Extend the horizon and check stability.} Compare historical prices/moments and terminal distribution distances. Market clearing can coexist with distortion from an endpoint imposed too early. The stationary continuation value does not reset population to stationarity.
\end{enumerate}

\lead{B. Calibrate: place an outer loop around all of A.}
Let $\theta$ contain the 11 free parameters, $\widehat M_m$ the empirical moments and $M_m$ the model measurements. The intended joint search minimizes:
\[
 \widehat\theta=\arg\min_{\theta\in\Theta}\sum_{m=1}^{12}
 w_m\left[M_m\bigl(p^*(\theta),\theta\bigr)-\widehat M_m\right]^2.
\]
Here $p^*(\theta)$ must be the equilibrium path for that candidate, not an inherited price guess. The weights include both empirical uncertainty estimates and declared synthetic scales.

\begin{enumerate}[itemsep=3pt]
\item \textbf{Propose $\theta$ within its bounds.} The 11 coordinates govern discounting; two fertility-shock scales; the owner-service premium; initial housing-supply scale; two bequest parameters; child housing needs; first-birth cost and housing jump; and the historical preference change. Tenure-shock scale and supply elasticity remain externally fixed.
\item \textbf{Renormalize and solve.} Refit the old preference intercept to completed fertility 2.1, rebuild the initial state and endpoint, and run all of algorithm A. Price-root iterations hold $\theta$ fixed; calibration iterations change $\theta$.
\item \textbf{Measure all 12 targets and score the candidate.} Four fertility moments cover completed fertility, childlessness and two first-birth timing statistics; five housing moments cover the birth response, family-size rooms gap, parent ownership gap, overall ownership and mean rooms; three wealth moments cover wealth/earnings, bequests/wealth and old-age dispersion.
\item \textbf{Use parameter sensitivities to propose the next $\theta$; return to step 2.} Assess rank and weak directions of the 12-by-11 equilibrium-adjusted moment Jacobian. This is different from the date-by-date \emph{price} Jacobian used in A. More targets than parameters does not by itself establish identification.
\end{enumerate}

\lead{Dates and current scope.} Targets use the 2024 CPS, 1979--84 birth cohorts, 2012--2023 ACS and approved wealth vintages. The birth-housing contrast maps to a 2019--2023 branch. Most model moments refer to 2023. Date/group alignment, horizon certification and new joint calibration remain pending.
\end{document}
